% sec-04: the foundation funnel and three mechanisms.
% Figure 9 (mermaid-type flowchart with gates), Table 13.
\section{The Foundation Route}

Two foundation letters go out on this day and neither asks for money. Both ask
whether a mechanism can fund a Phase 1 whose sponsor is a for-profit small
business with an academic clinical site under a separate agreement. That question
is answerable by email in a week; the same answer discovered after a full
application costs a cycle, and a cycle for a disease foundation is usually a
year.

\begin{appfloat}[!tb]
\begin{appfig}[mermaid-type / flowchart with labeled gates]
% Four gates along one row on a 3.05cm pitch.  Every exit leaves the row upward
% and every hold leaves it downward, so no edge crosses another edge and none
% passes through a node.
\node[mmin,text width=26mm] (q) at (0,0) {Mechanism question, asked before applying};
\node[mmdec,text width=20mm] (g1) at (3.05,0) {For-profit sponsor eligible?};
\node[mmdec,text width=20mm] (g2) at (6.10,0) {Phase 1 resectable in scope?};
\node[mmdec,text width=20mm] (g3) at (9.15,0) {Is a cycle open?};
\node[mmmid,text width=22mm]  (l)  at (12.20,0) {Letter of intent};
\node[mmdec,text width=20mm] (g4) at (15.25,0) {Named site at LOI stage?};
\node[mmgoal,text width=22mm] (a) at (18.30,0) {Full application, then award decision};
\node[mmgray,text width=24mm] (x1) at (3.05,1.75) {Stop. Effort saved, nothing filed};
\node[mmgray,text width=24mm] (x2) at (6.10,1.75) {Stop, or reshape to a different aim};
\node[mmsoft,text width=24mm] (w)  at (9.15,-1.80) {Hold the pack for the next cycle};
\node[mmsoft,text width=24mm] (s)  at (15.25,-1.80) {Hold until a site names an investigator};
\draw[mmedgeb] (q) -- (g1);
\draw[mmedgeb] (g1) -- node[mmlabel] {yes} (g2);
\draw[mmedgeb] (g2) -- node[mmlabel] {yes} (g3);
\draw[mmedgeb] (g3) -- node[mmlabel] {open} (l);
\draw[mmedgeb] (l) -- (g4);
\draw[mmedgeb] (g4) -- node[mmlabel] {no} (a);
\draw[mmedge] (g1) -- node[mmlabel] {no} (x1);
\draw[mmedge] (g2) -- node[mmlabel] {no} (x2);
\draw[mmedged] (g3) -- node[mmlabel] {closed} (w);
\draw[mmedged] (w) -| ([xshift=-0.45cm]l.south);
\draw[mmedged] (g4) -- node[mmlabel] {yes} (s);
\draw[mmedged] (s) -| ([xshift=-0.45cm]a.south);
\end{appfig}
\vspace{-0.60cm}
\figcaption{Figure 9. The foundation funnel drawn as four gates, with the exit at each and\\
the two holds, so that the cheapest gate to fail is visible as the first one.}
\end{appfloat}

\subsection{Three mechanisms, and what each would have to answer}

\begin{apptable}
\begin{tabularx}{\textwidth}{>{\raggedright\arraybackslash}p{3.4cm} >{\raggedright\arraybackslash}p{2.4cm} >{\raggedright\arraybackslash}X}
\headrow \thead{Mechanism class} & \thead{Cycle} & \thead{The eligibility question that decides it}\\
\midrule
Disease foundation research grant, pancreatic cancer specific & Published annual cycle & Whether the applicant of record may be a company with a named academic site, or must be the institution \\
Pre-competitive consortium, public and private partners & Rolling, by consortium formation & Whether a single-company applicant can join a mechanism designed for multi-party consortia \\
Institutional pilot award at a candidate site & Annual, applied for by the site & Whether a site investigator would apply, which is a site decision and not this company's \\
\end{tabularx}
\end{apptable}
\vspace{-0.60cm}
\tabcap{Table 13. Three foundation mechanism classes with their cycle shape and\\
the single eligibility question that determines whether an application fits.}

The third row is the one this company cannot act on alone. An institutional pilot
award is applied for by the institution, so it becomes available only after the
feasibility conversation of \S2 has produced an investigator who wants it. It is
listed here so that it is not forgotten at the moment it becomes possible
\cite{ucsdmoores2026}.

\subsection{What a foundation is told about the evidence}

The same six quantities that go to a federal reviewer, each carrying the
limitation its own authors stated: a ten-arm quantitative systems pharmacology
run at 12.8 months median overall survival against 5.4 \cite{qspmetpancre}, the
RASolute 302 readout at 13.2 against 6.6 \cite{rasolute302}, a 1000-patient
digital twin at 12.1 months and a progression-free hazard ratio of 0.31
\cite{onpremwhippl}, a credibility score of 81.9 across 55 verification tests
against the ASME V and V 40 framework \cite{asmevv40}, and an empirical
triplicate across 100,000 records \cite{daraxident}.

Three of the six are in silico or digital twin results and are labeled as such in
every table that carries them. The proximity between the 2025 simulation and the
2026 trial readout is a chronology observation and a hypothesis-supporting one,
not a validation claim, and the sample size, single-agent and molecular-selection
differences are material.
